\usepackage{eso-pic}
\makeatletter
\newsavebox{\papersbox}
\newdimen\ep@vx \newdimen\ep@vy \newdimen\ep@vxw \newdimen\ep@vyh
\ExplSyntaxOn
\fp_new:N \l_ep_scale_fp
\fp_new:N \l_ep_winW_fp
\fp_new:N \l_ep_winH_fp
\ExplSyntaxOff
\newcommand*{\@startpaper}{}%
\newcommand*{\@coverone}{}%
\newcommand*{\@covertwo}{}%
\newcommand*{\@coverthree}{}%
\newcount\l@covercount
\newcount\l@coveronepage
\newcount\l@covertwopage
\newcount\l@coverthreepage
\newcount\l@coverfourthpage
\newcommand*{\@coverfourth}{}%
% Una guarda por cada portada definida, en el orden en que se declaran:
% frontCover (1), backCover (2) y startpaper (3). La que no se define no
% ocupa página: las que sí, se encadenan desde la primera.
\newcommand{\setcover}[1]{%
  \advance\l@covercount by 1\relax%
  \titlepageguardstrue%
  \ifcase\l@covercount
    \relax
  \or
    \gdef\@coverone{#1}%
  \or
    \gdef\@covertwo{#1}%
  \or
    \gdef\@coverthree{#1}%
  \else
    \gdef\@coverfourth{#1}%
  \fi
}
\newcommand{\setstartpaper}[1]{%
  \gdef\@startpaper{#1}%
  \setcover{#1}%
}
% La página de la guarda que se acaba de definir: la actual. Se asigna
% cuando se crea la página, no cuando se define la guarda, porque entre una y
% otra puede haber páginas en blanco de por medio.
\newcommand{\coverassign}{%
  \ifcase\l@covercount
    \relax
  \or
    \global\l@coveronepage=\value{page}%
  \or
    \global\l@covertwopage=\value{page}%
  \or
    \global\l@coverthreepage=\value{page}%
  \else
    \global\l@coverfourthpage=\value{page}%
  \fi
}
% La página en blanco de una guarda. Se crea aquí y no en \maketitle para que
% siga existiendo aunque el proyecto redefina ese comando: si la página no
% existiera, el contenido se montaría encima de la imagen de la guarda.
\newcommand{\coverblankpage}{%
  \coverassign
  \thispagestyle{empty}%
  \null\clearpage
}
% El verso de la guarda. Con una sola guarda (frontCover o backCover) su verso
% queda en blanco: así el startpaper y lo que siguen caen en páginas impares.
% Con dos guardas no hace falta: el startpaper ya cae en un recto.
\newcommand{\coververso}{%
  \if@twoside
    \ifnum\l@covercount=1
      \thispagestyle{empty}%
      \null\clearpage
    \fi
  \fi
}
% La imagen se mide al dibujarla: así cada guarda usa su propia caja y todas
% comparten las mismas macros de dimensión.
\ExplSyntaxOn
\cs_new_protected:Npn \ep@measurecover #1 {
  \sbox \papersbox { \includegraphics {#1} }
  \fp_set:Nn \l_ep_scale_fp {
    max( ( \dim_to_fp:n { \the\paperwidth } ) / \dim_to_fp:n { \the\wd\papersbox },
         ( \dim_to_fp:n { \the\paperheight } ) / \dim_to_fp:n { \the\ht\papersbox } )
  }
  \fp_set:Nn \l_ep_winW_fp { ( \dim_to_fp:n { \the\paperwidth } ) / \l_ep_scale_fp }
  \fp_set:Nn \l_ep_winH_fp { ( \dim_to_fp:n { \the\paperheight } ) / \l_ep_scale_fp }
  \dim_set:Nn \ep@vx { \fp_to_dim:n { ( \dim_to_fp:n { \the\wd\papersbox } - \l_ep_winW_fp ) / 2 } }
  \dim_set:Nn \ep@vy { \fp_to_dim:n { ( \dim_to_fp:n { \the\ht\papersbox } - \l_ep_winH_fp ) / 2 } }
  \dim_set:Nn \ep@vxw { \fp_to_dim:n { ( \dim_to_fp:n { \the\wd\papersbox } + \l_ep_winW_fp ) / 2 } }
  \dim_set:Nn \ep@vyh { \fp_to_dim:n { ( \dim_to_fp:n { \the\ht\papersbox } + \l_ep_winH_fp ) / 2 } }
}
\ExplSyntaxOff
\newcommand{\drawcoverimage}[1]{%
  \ep@measurecover{#1}%
  \put(.5\paperwidth,\ifx\ESO@HookIIIBG\@empty .5\paperheight\else -.5\paperheight\fi){%
    \vbox to 0pt{%
      \vss
      \hbox to 0pt{%
        \hss
        \includegraphics[%
          viewport={\the\ep@vx} {\the\ep@vy} {\the\ep@vxw} {\the\ep@vyh},%
          clip,%
          width=\dimexpr\paperwidth\relax,%
          height=\dimexpr\paperheight\relax,%
        ]{#1}%
        \hss
      }%
      \vss
    }%
  }%
}
\newcommand{\drawcovers}{%
  \iftitlepageguards
    \ifx\@coverone\@empty\else
      \ifnum\value{page}=\l@coveronepage
        \drawcoverimage{\@coverone}%
      \fi
    \fi
    \ifx\@covertwo\@empty\else
      \ifnum\value{page}=\l@covertwopage
        \drawcoverimage{\@covertwo}%
      \fi
    \fi
    \ifx\@coverthree\@empty\else
      \ifnum\value{page}=\l@coverthreepage
        \drawcoverimage{\@coverthree}%
      \fi
    \fi
    \ifx\@coverfourth\@empty\else
      \ifnum\value{page}=\l@coverfourthpage
        \drawcoverimage{\@coverfourth}%
      \fi
    \fi
  \fi
}
\AddToShipoutPictureBG{\drawcovers}
% Membrete de las interventions: una imagen por página, al pie de la caja de
% texto y pegada al borde exterior (impar: derecha; par: izquierda). Si hay
% menos imágenes que páginas, la lista se repite desde el principio.
\newcounter{letterheadpage}
\newcounter{letterheadtotal}
\newcommand{\letterheadsetup}[1]{%
  \def\@letterheadlist{#1}%
  \setcounter{letterheadtotal}{0}%
  \@for\@letterheaditem:=\@letterheadlist\do{\stepcounter{letterheadtotal}}%
  \setcounter{letterheadpage}{0}%
  \AddToShipoutPictureFG{\@letterheadpage}%
}
\newcommand{\@letterheadpage}{%
  \stepcounter{letterheadpage}%
  \begingroup
    \@tempcnta=\value{letterheadpage}%
    \@whilenum\@tempcnta>\value{letterheadtotal}\do{\advance\@tempcnta-\value{letterheadtotal}}%
    \@tempcntb=\z@
    \def\@letterheadchosen{}%
    \@for\@letterheaditem:=\@letterheadlist\do{%
      \advance\@tempcntb\@ne
      \ifnum\@tempcntb=\@tempcnta \xdef\@letterheadchosen{\@letterheaditem}\fi
    }%
    \ifodd\value{page}%
      \AtTextLowerLeft{\put(\textwidth,0){\makebox[0pt][r]{\includegraphics[width=3cm]{\@letterheadchosen}}}}%
    \else
      \AtTextLowerLeft{\put(0,0){\makebox[0pt][l]{\includegraphics[width=3cm]{\@letterheadchosen}}}}%
    \fi
  \endgroup
}
\makeatother
